\section{Semantics}\label{sec:semantics}

A key contribution of this paper is a precise and neutral description of the operational behavior of several contemporary async systems. To reduce the cosmetic noise we provide the operational semantics as extensions to a base $\lambda$-calculus variant, henceforth refered to as \lc. By providing the semantics in layers we highlight the control features necessary for a given language's implementation and can more easily compare the operational behavior between two languages that have seemingly incomparable semantics. We additionally provide all the semantics in this section as a runnable Redex model, available in the supplementary materials.~\missingcite

\begin{figure}[H]
  \begin{subfigure}[t]{0.39\textwidth}
    {$$
\begin{array}{r@{\ }c@{\ }l}
\mathsf{Number}\ n & & \mathsf{String}\ s \\
\mathsf{Const}\ c      & ::=  & \mathsf{true} \mid\mathsf{false}\mid n \mid s \\
\mathsf{Expr}\ e       & ::=  & x \mid v \mid \mathbfsf{fn}\ \many{x} \to e \mid e\ \many{e} \\
                       & \mid & \mathbfsf{let}\ x = e\ \mathbfsf{in}\ e \\
                       & \mid & \mathbfsf{if}\ e\ \mathbfsf{then}\ e\ \mathbfsf{else}\ e \\
                       & \mid & \mathbfsf{fix}\ e \mid e \mathbin{;} e \\
                       & \mid & \mathbfsf{reset}\ e \mid \mathbfsf{shift}\ x \to e \\
                       & \mid & \mathbfsf{ref}\ e \mid !e \mid e := e \\
\\
\mathsf{Value}\ v      & ::=  & c \mid () \mid \mathsf{ptr}(x) \\
       & \mid & \mathbfsf{fn}\ \many{x} \to e \mid \mathbfsf{fix}\ v \\
\\
\mathsf{Eval\ Ctx}\ E      & ::=  & [\cdot] \mid \many{v}\ E\ \many{e} \mid \mathbfsf{fix}\ E \\
       & \mid & \mathbfsf{let}\ x = E\ \mathbfsf{in}\ e \\
       & \mid & \mathbfsf{let}\ x = v\ \mathbfsf{in}\ E \\
       & \mid & \mathbfsf{if}\ E\ \mathbfsf{then}\ e\ \mathbfsf{else}\ e \\
       & \mid & \many{v} \mathbin{;} E \mathbin{;} \many{e} \mid \mathbfsf{reset}\ E \\
       & \mid & \mathbfsf{ref}\ E \mid !E \mid E := e \\
       & \mid & v := E
\\
\mathsf{Delim\ Ctx}\ M      & ::=\  & \mathsf{(} E\ \mathsf{without}\ \mathbfsf{reset}\ M \mathsf{)}
\\
\mathsf{Heap}\ \sigma & ::=  & \cdot \mid \sigma, x \mapsto v
\end{array}
$$

     \cprotect\caption{The grammar for \lc, a call-by-value variant of the $\lambda$-calculus with mutable references and delimited continuations.}
     \label{fig:lang:lc}}
  \end{subfigure}\hfill%
  \begin{subfigure}[t]{0.59\textwidth}
    {$$
\begin{array}{r@{\ }c@{\ }l}
  e & ::= \ldots \mid \bsfid{throw-in}\ e\ e\ \mid \bsfid{throw}\ e \mid \bsfid{try}\ e\ \bsfid{catch}\ e \\
  % \mathsf{Exn\ Ctx}\ G & ::=\  \mathsf{E} & \mid \bsfid{throw-in}\ e\ e\ \mid \bsfid{throw}\ e
\end{array}
$$

     \cprotect\subcaption{The grammar for \exn, which extends \lc with \code|try|/\code|catch|, and \code|throw| forms.}
    \label{fig:lang:exn}}
    \vspace{0.5em}
    {\noindent
\begin{tabular*}{\linewidth}{@{} l @{\extracolsep{\fill}} r @{}}
  $\langle \sigma, E[\bsfid{throw-in}(\mathbf{fn}\ x \to E_{\mathit{inner}}[x], v)] \rangle$ &  [Throw-In] \\
$\longrightarrow \langle \sigma, E[E_{\mathit{inner}}[\bsfid{throw}\ v]] \rangle$ &
\end{tabular*}

     \cprotect\subcaption{The operational semantics for the \exn.}
     \label{fig:lang:exn:red}}
  \end{subfigure}

  \caption{The syntax and semantics for the core \lc and \exn calculi.}
  \label{fig:lang:core}
\end{figure}

\Cref{fig:lang:core} provides the grammar and semantics of the core calculi that comprise our base for the async extensions. The core language \lc contains mutable references and delimited continuations via the \code|shift| and \code|reset| operators described by \citet{danvy_functional_nodate}. The grammar for \lc is provided in \Cref{fig:lang:lc}. The reduction rule is defined over a mutable heap $\sigma$ and an expression $e$, $\langle\sigma\,e\rangle$. The evaluation context $M$ is the continuation delimited by an enclosing \code|reset| form and captured by \code|shift|. For space reasons, the standard reduction rules and full grammar for \lc are available in \Cref{appendix}.

For languages with exceptional control flow, we extend \lc with a \code|try|/\code|catch| and \code|throw| form --- we will refer to this language as \exn. The grammar and reduction rules for \exn are provided in \Cref{fig:lang:exn}. The main contribution of this grammar is the evaluation context $G$, which is the exception context delimited by a \code|try| form and discarded by the \code|throw| operator. We also provide a \code|throw_in| expression that raises an exception within a suspended computation. We model this rule based on Python's \code{.throw()} coroutine method, which resumes a suspended computation with an exception instead of a value.

We now describe the configuration of the shared async grammar. \Cref{fig:lasync:grammar} defines the shared grammar constructs that encode the async runtime building blocks discussed in \Cref{sec:designdims}. We will refer to this shared grammar extension as the async platform, \lasync.

\begin{figure}[H]
  \begin{subfigure}[t]{\textwidth}
    $$
\begin{array}{l}
\begin{array}{rcl @{\hspace{3em}} rcl}
  \sfid{Time}\ t              & \in & \mathbb{N}_0 \cup \{\infty\}                              & \sfid{Frame}\ F         & ::= & \langle \ell, e \rangle \\
\sfid{Label}\ \ell          & ::= & x \mid \bsfid{sync}                    & \sfid{Frame\ Stack}\ FS & ::= & \epsilon \mid F \circ FS \\
\sfid{Continuation}\ \kappa & ::= & \bsfid{fn}\ x_{\mathit{unit}} \to e    & \sfid{Process}\ P       & ::= & \{ FS_1, \dots, FS_n \} \\
\sfid{Signal\ Queue}\ T     & ::= & \epsilon \mid (t, \ell, \kappa) \circ T & \sfid{Task\ Queue}\ Q   & ::= & \epsilon \mid (\ell, \kappa) \circ Q 
\end{array} \\
\\
\begin{array}{rcl}
e & ::= & \dots \mid \bsfid{sys-io}(e, e) \mid \bsfid{sys-start-soon}(\many{\ell,e}) \mid \bsfid{sys-start-later}(e, \ell, e)
% e & ::= & \dots \mid \bsfid{sys-block}(e) \mid \mathbfsf{sys-time}() \mid \mathbfsf{sys-io}(e, e) \mid \mathbfsf{sys-start-soon}(\many{\ell,e}) \mid \mathbfsf{sys-start-later}(e, \ell, e)
\end{array}
\end{array}
$$

    \cprotect\caption{The full grammar for \lasync.}
    \label{fig:lasync:grammar}
  \end{subfigure}\vspace{1em}
  \begin{subfigure}[t]{\textwidth}
    \begin{tabular*}{\linewidth}{@{} l @{\extracolsep{\fill}} r @{}}
  $\langle t, \sigma_0, Q, T, P \uplus \{ \langle \ell, E[\text{os\_io}(t_{\mathit{delay}}, v)] \rangle \circ FS \} \rangle$ & [OS-IO] \\
  $\longrightarrow \langle \text{step}(t), \sigma_1, Q, T, P \uplus \{ F_{\mathit{new}} \circ \langle \ell, E[x_{\mathit{io}}] \rangle \circ FS \} \rangle$ & \\
\multicolumn{2}{@{}l@{}}{
  $\begin{array}{l}
  \textbf{where } (\sigma_1, x_{io}) = \text{task\_alloc}(\sigma_0) \\
  \qquad e_{try} = \mathbfsf{try}\ none \mathbin{;} \text{task\_set\_done}(x_{io}, v)\ \mathbfsf{catch}\ e_{err} \to \text{task\_set\_failed}(x_{io}, e_{err}) \\
  \qquad e_{done} = e_{\mathit{try}} \mathbin{;} \text{os\_start\_soon}(\text{task\_get\_dependents}(x_{io})) \\
  \qquad F_{new} = \langle x_{\mathit{io}}, \text{os\_start\_later}(\text{os\_time}() + t_{\mathit{delay}}, x_{\mathit{io}}, \text{fun } () \to e_{done}) \rangle
  \end{array}$
}
\end{tabular*}

    \cprotect\caption{An example of the reduction rules for \lasync.}
    \label{fig:lasync:red}
  \end{subfigure}
  \caption{Syntax and state components for the Async Platform Language extension.}
  \label{fig:lasync}
\end{figure}

A \textit{frame}, $F$, is written $\langle \ell,\,e\rangle$ and consists of an expression $e$, along with a frame label, $\ell$. A \textit{frame label}, $\ell$, is either the name of the \task this frame corresponds to, or \textit{sync} to denote a synchronous frame. A \textit{frame stack}, $FS$, is a list of frames --- essentially a thread. An empty frame stack is written $\epsilon$, and we write $F \circ FS$ to denote a frame stack whose head is $F$ and tail is $FS$. A \textit{process}, $P$, is a collection of frame stacks. We will use the syntax $P \uplus FS$ to select one frame stack in $P$ to run. A \textit{task queue}, $Q$, stores the async work waiting to be run; $Q$ comprises tuples that have a label and a continuation, $(\ell,\kappa)$. A \textit{signals queue}, $T$, stores async work that is waiting on \abbr{io} operations, the work will be rescheduled in the future at timestampt $t$, at the earliest. $T$ comprises tuples that have a time, a label, and a continuation, $(t,\ell,\kappa)$. We will use the same stack notation for $Q$ and $T$ as we do for $FS$, but will write $Q \qpush (\ell,\,\kappa)$, and similarly for $T$, to append to the end of the queue.

The async platform provides a few additional expressions that we will describe here. The $\text{os\_block}\,e$ expression provides the barrier between synchronous code and asynchronous code. Languages may use a block expression to block a synchronous frame until the asynchronous expression has completed running. The block expression marks the runtime entry, and when it returns, the runtime exit.

The expression \code|os_io|$(e_t,e)$ is the \textsc{OmniScient io} operation provided by the runtime system. It models real-world \abbr{io} operations that would take some non-deterministic amount of time to resolve, like receiving packets from the network or querying a database. The form is omniscient because the resulting expression is provided upfront, e.g., \code|os_io|$(5, \text{``hello''})$ creates a \task that will resolve to the value ``hello'' after at least five time steps. Thus the form is deterministic in its evaluation \textit{value}, but not in its evaluation \textit{time.}

The expression \code|os_start_soon|$(\ell,\kappa)$ appends the continuation $\kappa$ to $Q$, set to run in a frame with the specified label $\ell$. The expression \code|os_start_later|$(t,\ell,\kappa)$ appends the continuation to the signal queue, $T$, which will be scheduled after timestep $t$. The expression \code|os_time|$()$ simply evaluates to the current timestep. 

The reduction rules for \lasync, and all languages that extent it, are defined over the tuple $\langle t, \sigma, Q, T, P \rangle$. The goal of the operational semantics is to highlight the differences between languages. As such we will present the same reduction rule from different languages together. We will also call attention to the portion of the redcution relevant to each design dimension with a highlight in the following colors: \ceag{\eagerness,} \cext{\extent,} \crfe{\references,} \cdes{\destruction,} \cpro{\propagation,} \csus{\suspension,} and \ccan{\cancellation.} We simply the presentation with a liberal use of metafunction; the full definitions of grammars, reductions, and metafunctions are provided in \Cref{app:semantics}.

\subsection{Async Function Application}

\begin{figure}[H]
  \noindent
$\begin{array}[t]{@{}r@{\ }l@{}l@{}l@{}l@{}}
                & \elided{\sigma_0, P \uplus \{ & & \langle \ell, E[(\bsfid{async\ fn}\ \many{x} \to e_{\mathit{body}})\ \many{v}] \rangle & {}\circ FS \}} \\
\longrightarrow & \elided{\sigma_2, P \uplus \{ \ceag{\langle x_{\mathit{task}}, e_{\mathit{task}} \rangle} & {}\circ{} & \langle \ell, E[x_{\mathit{task}}] \rangle & {}\circ FS \}}
\end{array}$\hfill$\text{[\CSharp/\Js-Async-App]}$\par
\begin{tabular*}{\linewidth}{@{} l @{\extracolsep{\fill}} r @{}}
\multicolumn{2}{@{}l@{}}{
  $\begin{array}{l}
  \where\ \many{x}_{\mathit{fresh}} = \many{\sfid{gensym}(),~ \sigma_1 = \sigma_0[\many{x}_{\mathit{fresh}} \mapsto \many{v}],~ e_{\mathit{subst}} = e_{\mathit{body}}[\many{x}_{\mathit{fresh}}/\many{x}]} \\
  \qquad (\sigma_2, x_{\mathit{task}}) = \sfid{task-alloc}(\sigma_1)\;\cext{\ghost} \\
  \qquad e_{\mathit{task}} = \bsfid{reset}\ ( \\
  \qquad \qquad \bsfid{try}\ \semmacro{task-set-done}(x_{\mathit{task}}, e_{\mathit{subst}})\ \bsfid{catch}\ v_{\mathit{err}} \to \semmacro{task-set-failed}(x_{\mathit{task}}, v_{\mathit{err}}) \\
  \quad\;\;\cdes{\ghost}\;\cext{\ghost}\;\cpro{\ghost}\;
  \bsfid{sys-start-soon}(\semmacro{task-get-dependents}(x_{\mathit{task}})))
  \end{array}$
}
\end{tabular*}

\vspace{1em}
\noindent
$\begin{array}[t]{@{}r@{\ }l@{\ }l@{}l@{}}
                & \elided{\sigma_0, Q, & P \uplus \{ \langle \ell, E[(\bsfid{async\ fn}\ \many{x} \to e_{\mathit{body}})\ \many{v}] \rangle & {}\circ FS \}} \\
\longrightarrow & \elided{\sigma_3, \ceag{Q \qpush (x_{\mathit{task}}, \kappa)}, & P \uplus \{ \langle \ell, E[x_{\mathit{task}}] \rangle & {}\circ FS \}}
\end{array}$\hfill$\text{[Swift-Async-App]}$\par
\begin{tabular*}{\linewidth}{@{} l @{\extracolsep{\fill}} r @{}}
\multicolumn{2}{@{}l@{}}{
  $\begin{array}{l}
  \where\ \many{x}_{\mathit{fresh}} = \many{\sfid{gensym}(),~ \sigma_1 = \sigma_0[\many{x}_{\mathit{fresh}} \mapsto \many{v}],~ e_{\mathit{subst}} = e_{\mathit{body}}[\many{x}_{\mathit{fresh}}/\many{x}]} \\
  \qquad (\sigma_2, x_{\mathit{task}}) = \sfid{task-alloc}(\sigma_1) \\
  \qquad \cext{\sigma_3 = \sfid{task-add-dependency}(\sigma_2, \ell, x_{\mathit{task}})} \\
  \qquad e_{\mathit{task}} = \bsfid{reset}\ ( \\
  \qquad \qquad \bsfid{try}\ \semmacro{task-set-done}(x_{\mathit{task}}, e_{\mathit{subst}})\ \bsfid{catch}\ v_{\mathit{err}} \to \semmacro{task-set-failed}(x_{\mathit{task}}, v_{\mathit{err}}) \mathbin{;} \\
  \qquad \qquad \cdes{\semmacro{task-cancel-dependencies}(x_{\mathit{task}}) \mathbin{;}} \\
  \qquad \qquad \cext{\semmacro{task-wait-on-dependencies}(x_{\mathit{task}}) \mathbin{;}}\;\cpro{\ghost}\\
  \qquad \qquad \bsfid{sys-start-soon}(\semmacro{task-get-dependents}(x_{\mathit{task}}))) \\
  \qquad \kappa = \bsfid{fn}\ x_\mathit{unit} \to x_\mathit{unit}\mathbin{;}\ e_{\mathit{task}}
  \end{array}$
}
\end{tabular*}

\vspace{1em}
\noindent
$\begin{array}[t]{@{}r@{\ }l@{}}
                & \langle \sigma_0, E[(\bsfid{async\ fn}\ \many{x} \to e_{\mathit{body}})\ \many{v}] \rangle \\
\longrightarrow & \langle \sigma_1, E[\ceag{\bsfid{reset}\ (\bsfid{shift}\ k \to k \mathbin{;} e_{\mathit{subst}})}] \rangle
\end{array}$\hfill$\text{[Rust/Python-Async-App]}$\par
\begin{tabular*}{\linewidth}{@{} l @{\extracolsep{\fill}} r @{}}
\multicolumn{2}{@{}l@{}}{
  $\begin{array}{l}
  \where\ \many{x}_{\mathit{fresh}} = \many{\sfid{gensym}(),~ \sigma_1 = \sigma_0[\many{x}_{\mathit{fresh}} \mapsto \many{v}],~ e_{\mathit{subst}} = e_{\mathit{body}}[\many{x}_{\mathit{fresh}}/\many{x}]}
  \end{array}$
}
\end{tabular*}

  \makelegend{eag,ext,ref,des,pro,sus,can}
  \cprotect\caption{A comparison of the rule \semrule{Async-App} across languages.}
  \label{fig:async-app}
\end{figure}

\Cref{fig:async-app} shows the \semrule{Async-App} rule for \CSharp, Swift, and Python; the rule primarily has implications for \eagerness. In \CSharp a new \task is allocated, the function body, which is placed inside error handling logic to trap exceptions, is set in a new frame on the \emph{top} of the current frame stack. After the entire body has run and a value set in the \task, the dependents of the \task are rescheduled in the task queue. Note that dependencies of the \task are not waited on, cancelled, or terminated as they have indefinite extent.

The \semrule{Async-App} rule for Swift follows a similar procedure as that of \CSharp's, but with several important changes. First, the \task allocated is not independent, it is a \emph{dependency} of the currently executing \task labelled $\ell$. Registering spawned \tasks as dependencies is important because at the end of execution a \task will cancel and wait on it's dependencies before rescheduling its dependents in the task queue. Note that dependencies are \emph{waited on} and not \emph{awaited,} by merely waiting on dependencies to finish, any errors that occur are not propagated. The second key difference is that the newly created \task is appended to the task queue, not placed on the current frame stack.

\begin{figure}[H]
  \noindent
$\begin{array}[t]{@{}r@{\ }l@{\ }l@{}l@{}}
                & \elided{ \sigma_0, Q, & P \uplus \{ \langle \ell, E[\bsfid{spawn}\ v_{\mathit{coro}}] \rangle & {}\circ FS \}} \\
\longrightarrow & \elided{ \sigma_1, \ceag{Q \qpush (x_{\mathit{task}}, \kappa)}, & P \uplus \{ \langle \ell, E[x_{\mathit{task}}] \rangle & {}\circ FS \}}
\end{array}$\hfill$\text{[Asyncio/Tokio/Smol-Spawn]}$\par
\begin{tabular*}{\linewidth}{@{} l @{\extracolsep{\fill}} r @{}}
\multicolumn{2}{@{}l@{}}{
  $\begin{array}{l}
  \where\ (\sigma_1, x_{\mathit{task}}) = \sfid{task-alloc}(\sigma_0) \\
  \qquad e_{\mathit{task}} = \bsfid{reset}\ ( \\
  \qquad \qquad \difff{\bsfid{try}}\ \semmacro{task-set-done}(x_{\mathit{task}}, \bsfid{await}\ v_{\mathit{coro}})\ \difff{\bsfid{catch}\ v_{\mathit{err}} \to \semmacro{task-set-failed}(x_{\mathit{task}}, v_{\mathit{err}})} \mathbin{;} \\
  \quad\;\;\cdes{\ghost}\;\cext{\ghost}\;\cpro{\ghost}\;
  \bsfid{sys-start-soon}(\semmacro{task-get-dependents}(x_{\mathit{task}}))) \\
  \qquad \kappa = \bsfid{fn}\ x_{\mathit{unit}} \to x_{\mathit{unit}} \mathbin{;} e_{\mathit{task}}
  \end{array}$
}
\end{tabular*}

\vspace{1em}
\noindent
$\begin{array}[t]{@{}r@{\ }l@{\ }l@{}l@{}}
                & \elided{ \sigma_0, Q, & P \uplus \{ \langle \ell, E[\bsfid{spawn}\ v_{\mathit{coro}}] \rangle & {}\circ FS \}} \\
\longrightarrow & \elided{ \sigma_2, \ceag{Q \qpush (x_{\mathit{task}}, \kappa)}, & P \uplus \{ \langle \ell, E[x_{\mathit{task}}] \rangle & {}\circ FS \}}
\end{array}$\hfill$\text{[Trio-Spawn]}$\par
\begin{tabular*}{\linewidth}{@{} l @{\extracolsep{\fill}} r @{}}
\multicolumn{2}{@{}l@{}}{
  $\begin{array}{l}
  \where\ (\sigma_1, x_{\mathit{task}}) = \sfid{task-alloc}(\sigma_0) \\
    \qquad \cext{\sigma_2 = \sfid{task-add-dependency}(\sigma_1, \ell, x_{\mathit{task}})} \\
  \qquad e_{\mathit{task}} = \bsfid{reset}\ ( \\
  \qquad \qquad \bsfid{try}\ \semmacro{task-set-done}(x_{\mathit{task}}, \bsfid{await}\ v_{\mathit{coro}})\ \bsfid{catch}\ v_{\mathit{err}} \to \semmacro{task-set-failed}(x_{\mathit{task}}, v_{\mathit{err}}) \mathbin{;} \\
  \qquad \quad\;\cdes{\ghost}\;\cext{\semmacro{task-wait-on-dependencies}(x_{\mathit{task}}) \mathbin{;}}\\
  \qquad \qquad \cpro{\semmacro{task-reraise-dependency-failures}(x_{\mathit{task}}) \mathbin{;}} \\
  \qquad \qquad \bsfid{sys-start-soon}(\semmacro{task-get-dependents}(x_{\mathit{task}}))) \\
  \qquad \kappa = \bsfid{fn}\ x_{\mathit{unit}} \to x_{\mathit{unit}} \mathbin{;} e_{\mathit{task}}
  \end{array}$
}
\end{tabular*}

  \makelegend{eag,ext,ref,des,pro,sus,can}
  \cprotect\caption{A comparison of \semrule{Spawn} across Python runtimes.}
  \label{fig:spawn}
\end{figure}

The rule \semrule{Async-App} for Python is compared to the others very simple. Python's async design separates the \coro mechanism --- generated by the compiler from an \psudo{async} function --- and the runtime. Because Python exposes \coros as a building block for runtime libraries, an async function application evaluates immediately to a \coro. To compare the lifecycles of \tasks between languages the rules \semrule{Spawn} for Asyncio and Trio are included in \Cref{spawn}. An important detail in the definitions for \semrule{Spawn} is Trio's enforcement of \dyn \extent. In Trio, when a \task finishes executing it will \emph{await} its dependencies without cancelling them.

\ggnote{TODO PICK UP HERE: I caved and decided that SK's suggestion for presenting the semantics was better}

\subsection{\CSharp}

\begin{figure}[H]
  \begin{subfigure}{\textwidth}
    $$
\begin{array}{rcl @{\hspace{3em}} rcl}
  e & ::= & \ldots \mid \mathbfsf{async\ fn}\ \many{x} \to e \mid \mathbfsf{await}\ e & v & ::= & \ldots \mid \mathbfsf{async\ fn}\ \many{x} \to e
\end{array}
$$

    \cprotect\caption{The grammar extension for \CSharp's async/await expressions.}
    \label{fig:cs:grammar}
  \end{subfigure}\vspace{1em}
  \begin{subfigure}{\textwidth}
    \noindent
\begin{tabular*}{\linewidth}{@{} l @{\extracolsep{\fill}} r @{}}
$\langle t, \sigma_0, Q, T, P \uplus \{ \langle \ell, E[(\mathbf{async\ fn}\ \many{x} \to e_{\mathit{body}})\ \many{v}] \rangle \circ FS \} \rangle$ & [Async-App] \\
$\longrightarrow \langle \mathrm{step}(t), \sigma_2, Q, T, P \uplus \{ \ceag{F_{\mathit{new}}} \circ \langle \ell, E[x_{\mathit{task}}] \rangle \circ FS \} \rangle$ & \\
\multicolumn{2}{@{}l@{}}{
  $\begin{array}{l}
  \textbf{where}\\
  \qquad (\sigma_1, x_{\mathit{task}}) = \mathrm{task\_alloc}(\sigma_0)\ \mathrm{and}\ \many{x}_{\mathit{fresh}} = \many{\mathrm{gensym}()} \\
  \qquad \sigma_2 = \sigma_1[\many{x}_{\mathit{fresh}} \mapsto \many{v}] \\
  \qquad e_{\mathit{subst}} = e_{\mathit{body}}[\many{x}_{\mathit{fresh}}/\many{x}] \\
  \qquad e_{\mathit{try}} = \mathbfsf{try}\ \mathrm{task\_set\_done}(x_{\mathit{task}}, e_{\mathit{subst}})\ \mathbfsf{catch}\ v_{\mathit{err}} \to \mathrm{task\_set\_failed}(x_{\mathit{task}}, v_{\mathit{err}}) \\
  \qquad F_{\mathit{new}} = \langle x_{\mathit{task}}, \mathrm{reset}\ (e_{\mathit{try}} \mathbin{;} \cext{\mathrm{os\_start\_soon}(\mathrm{task\_get\_dependents}(x_{\mathit{task}}))}) \rangle
  \end{array}$
}
\end{tabular*}

\vspace{1.5em}

\noindent
\begin{tabular*}{\linewidth}{@{} l @{\extracolsep{\fill}} r @{}}
$\langle t, \sigma, Q, T, P \uplus \{ \langle \ell, E[\mathbf{await}\ v_{\mathit{aw}}] \rangle \circ FS \} \rangle$ & [Await] \\
$\longrightarrow \langle \mathrm{step}(t), \sigma, Q, T, P \uplus \{ \langle \ell, E[e_{\mathit{check}}] \rangle \circ FS \} \rangle$ & \\
\multicolumn{2}{@{}l@{}}{
  $\begin{array}{l}
    \textbf{where}\\
   \qquad e_k = \mathrm{shift}\ k \to \mathrm{task\_add\_self\_as\_dependent}(v_{\mathit{aw}}, \ell,\mathrm{task\_continue\_with}(v_{\mathit{aw}},k)) \\
  \qquad e_{\mathit{check}} = \csus{\mathrm{if}\ \mathrm{task\_is\_completed}(v_{\mathit{aw}})\ \mathrm{then}\ \mathrm{task\_get\_result}(v_{\mathit{aw}})\ \mathrm{else}\ e_k}
  \end{array}$
}
\end{tabular*}

\vspace{1.5em}

\noindent
\begin{tabular*}{\linewidth}{@{} l @{\extracolsep{\fill}} r @{}}
$\langle t, \sigma, Q, T, P \uplus \{ \langle \mathrm{root}, E[\mathrm{os\_block}\ v_{\mathit{aw}}] \rangle \circ FS \} \rangle$ & [OS-Block-Done] \\
$\longrightarrow \langle \mathrm{step}(t), \sigma, \cdes{\epsilon}, \cdes{\epsilon}, \{ \langle \mathrm{root}, E[\mathrm{task\_get\_result}(v_{\mathit{aw}})] \rangle \circ \cdes{\epsilon} \} \rangle$ & \\
\multicolumn{2}{@{}l@{}}{
  $\begin{array}{l}
  \textbf{where}\ \mathrm{task\_settled}(\sigma, v_{\mathit{aw}}) \\
  \end{array}$
}
\end{tabular*}

    \makelegend{eag,ext,des,sus}
    \cprotect\caption{The operation semantics for \CSharp's async/await expressions.}
    \label{fig:cs:red}
  \end{subfigure}
  \caption{The syntax and operational semantics for async/await \CSharp.}
  \label{fig:lang:cs}
\end{figure}


There are two lines that compose Swifts enforcement of \dyn \extent. First, the \task allocated is a \emph{dependency} of the currently executing frame labeled $\ell$. Second, when the application is finished the new \task, $x_{\mathit{task}}$, awaits all of its dependencies before returning and rescheduling it's dependents. Furthermore, before awaiting its dependencies the \task will cancel it's dependencies as a way to model \cancelled \destruction. Although the dependencies are awaited, if an unawaited exception was raised, it will not propagate at the end of the scope. After awaiting dependencies the current \task immediately reschedules dependents to run.

The langauge extension and reduction rules for \CSharp are provided in \Cref{fig:lang:cs}. The \semrule{Async-App} rule demonstrates \CSharp's eager evaluation. When an async function is applied, the body is placed in a new frame, which is placed at the top of the current frame stack. By placing the function body at the top of the current frame stack, evaluation of the current thread will continue with the callee. Notably, when the new \task is created, labelled as  $x_{\mathit{task}}$, it does not register itself as a dependency of the executing \task context, labeled $\ell$. Following a \task's completion, the dependents of a task are scheduled to resume and the \task frame exits. The absence of dependency registration and the absence of awaiting on dependencies before waking dependents compose \CSharp's \indef \extent.

The \semrule{Await} rule demonstrates \CSharp's dynamic \suspension policy. When executing an await expression \CSharp will check if the \task\cprotect\footnote{Any \code{Awaitable} object can be awaited in \CSharp. \citet{hutchison_pause_2012} demonstrate how to encode this pattern in featherweight \CSharp$_5$, an extension to our model is straightforward and left as an exercise for the reader.} is completed with a value, if it is, then the value is retrieved and execution continues.

The \semrule{OS-Block-Done} shows \terminated \destruction semantics. When the \task blocking the synchronous \code|root| frame settles, i.e., has completed with a value or failed with an exception, scheduled computations in $Q$, pending \abbr{io} operations in $T$, and running compuations on other threads are all terminated.

\subsection{Swift}

\begin{figure}[H]
  \begin{subfigure}{\textwidth}
    $$
\begin{array}{rcl @{\hspace{3em}} rcl}
  e & ::= & \ldots \mid \mathbf{async}\ \mathbf{fn}\ \many{x} \to e \mid \mathbf{await}\ e \mid \mathrm{cancel}\ e \mid \mathrm{cancelled?}
  & v & ::= & \ldots \mid \mathbf{async}\ \mathbf{fn}\ \many{x} \to e
\end{array}
$$

    \cprotect\caption{The Swift grammar extension on \exn.}
    \label{fig:swift:grammar}
  \end{subfigure}\vspace{1em}
  \begin{subfigure}{\textwidth}
    \noindent
\begin{tabular*}{\linewidth}{@{} l @{\extracolsep{\fill}} r @{}}
$\langle t, \sigma_0, Q_0, T, P \uplus \{ \langle \ell, E[(\mathbf{async\ fn}\ \many{x} \to e_{\mathit{body}})\ \many{v}] \rangle \circ FS \} \rangle$ & [Async-App] \\
$\longrightarrow \langle \mathrm{step}(t), \sigma_2, Q_1, T, P \uplus \{ \langle \ell, E[x_{\mathit{task}}] \rangle \circ FS \} \rangle$ & \\
\multicolumn{2}{@{}l@{}}{
  $\begin{array}{l}
  \textbf{where}\ (\sigma_1, x_{\mathit{task}}) = \cext{\mathrm{task\text{-}alloc\text{-}dependency}(\sigma_0,\ell)}\ \mathrm{and}\ \many{x}_{\mathit{fresh}} = \many{\mathrm{gensym}()} \\
  \qquad \sigma_2 = \sigma_1[\many{x}_{\mathit{fresh}} \mapsto \many{v}] \\
  \qquad e_{\mathit{subst}} = e_{\mathit{body}}[\many{x}_{\mathit{fresh}}/\many{x}] \\
  \qquad \kappa = \mathbf{fn}\ x_\mathit{unit} \to x_\mathit{unit}\mathbin{;}\ \mathbfsf{reset}\ ( \\
  % \qquad \qquad \mathrm{task\text{-}add\text{-}parent}(x_{\mathit{task}}, \ell) \mathbin{;} \\
  \qquad \qquad \mathbfsf{try}\ \mathrm{task\text{-}set\text{-}done}(x_{\mathit{task}}, e_{\mathit{subst}})\ \mathbfsf{catch}\ v_{\mathit{err}} \to \mathrm{task\text{-}set\text{-}failed}(x_{\mathit{task}}, v_{\mathit{err}}) \mathbin{;} \\
  \qquad \qquad \cdes{\mathrm{task\text{-}cancel\text{-}dependencies}(x_{\mathit{task}}) \mathbin{;}} \\
  \qquad \qquad \cpro{\mathrm{task\text{-}wait\text{-}on\text{-}dependencies}}\cext{(x_{\mathit{task}}) \mathbin{;}}\\
  \qquad \qquad \mathrm{os\text{-}start\text{-}soon}(\mathrm{task\text{-}get\text{-}dependents}(x_{\mathit{task}})) \\
  \qquad ) \\
  \qquad Q_1 = \ceag{Q_0 \qpush (x_{\mathit{task}}, \kappa)}
  \end{array}$
}
\end{tabular*}

% \vspace{1em}
% \noindent
% \begin{tabular*}{\linewidth}{@{} l @{\extracolsep{\fill}} r @{}}
% $\langle t, \sigma, Q, T, P \uplus \{ \langle \ell, E[\mathrm{await}\ v_{\mathit{aw}}] \rangle \circ FS \} \rangle$ & [Await] \\
% $\longrightarrow \langle \mathrm{step}(t), \sigma, Q, T, P \uplus \{ \langle \ell, E[e_{\mathit{check}}] \rangle \circ FS \} \rangle$ & \\
% \multicolumn{2}{@{}l@{}}{
%   $\begin{array}{l}
%   \textbf{where}\ e_k = \mathbfsf{shift}\ k \to \mathrm{task\text{-}add\text{-}waiter}(v_{\mathit{aw}}, \ell, \mathrm{task\text{-}continue\text{-}with}(v_{\mathit{aw}}, k)) \\
%   \qquad e_{\mathit{check}} = \csus{\mathbfsf{if}\ \mathrm{task\text{-}is\text{-}completed}(v_{\mathit{aw}})\ \mathbfsf{then}\ \mathrm{task\text{-}get\text{-}result}(v_{\mathit{aw}})\ \mathbfsf{else}\ e_k}
%   \end{array}$
% }
% \end{tabular*}

\vspace{1em}
% \noindent
% \begin{tabular*}{\linewidth}{@{} l @{\extracolsep{\fill}} r @{}}
% $\langle t, \sigma, Q, T, P \uplus \{ \langle \ell, E[\mathrm{cancel}\ v_{\mathit{task}}] \rangle \circ FS \} \rangle$ & [Cancel] \\
% $\longrightarrow \langle t, \sigma, Q, T, P \uplus \{ \langle \ell, E[\mathrm{task\text{-}set\text{-}cancelled}(v_{\mathit{task}})] \rangle \circ FS \} \rangle$ & 
% \end{tabular*}
\noindent
\begin{tabular*}{\linewidth}{@{} l @{\extracolsep{\fill}} r @{}}
$\langle t, \sigma, Q, T, P \uplus \{ \langle \ell, E[\mathrm{cancel}\ v_{\mathit{task}}] \rangle \circ FS \} \rangle$ & [Cancel] \\
$\longrightarrow \langle t, \sigma, Q, T, P \uplus \{ \langle \ell, E[v_{\mathit{task}}.\mathrm{cancelled} := \mathrm{true}] \rangle \circ FS \} \rangle$ & 
\end{tabular*}

\vspace{1em}
\noindent
\begin{tabular*}{\linewidth}{@{} l @{\extracolsep{\fill}} r @{}}
$\langle t, \sigma, Q, T, P \uplus \{ \langle \ell, E[\mathrm{cancelled}()] \rangle \circ FS \} \rangle$ & [Cancelled?] \\
$\longrightarrow \langle t, \sigma, Q, T, P \uplus \{ \langle \ell, E[\ccan{\mathrm{task\text{-}cancelled}(\sigma, \ell)}] \rangle \circ FS \} \rangle$ & 
\end{tabular*}

\vspace{1em}

\noindent
\begin{tabular*}{\linewidth}{@{} l @{\extracolsep{\fill}} r @{}}
$\mathrm{task\text{-}cancelled}(\sigma, x) = \mathit{self\text{-}cancelled} \lor \mathit{dependent\text{-}cancelled}$ & \\
\multicolumn{2}{@{}l@{}}{
  $\begin{array}{l}
  \textbf{where}\\
  \qquad \mathit{self\text{-}cancelled} =\; !\sigma(x).\mathrm{cancelled} \\
  \qquad \many{x}_{\mathit{parents}} = \sigma(v_{\mathit{task}}.\mathrm{parents}) \\
  \qquad \mathit{dependent\text{-}cancelled} = \bigvee_{x_p \in \many{x}_{\mathit{parents}}} \mathrm{task\text{-}cancelled}(\sigma, x_p)
  \end{array}$
}
\end{tabular*}

    \makelegend{eag,ext,des,pro,sus,can}
    \cprotect\caption{}
    \label{fig:swift:red}
  \end{subfigure}
  \cprotect\caption{The core operational semantics for Swift async/await.}
  \label{fig:lang:swift}
\end{figure}

The langauge extension and reduction rules for Swift are provided as an extension of \exn in \Cref{fig:lang:swift}. Compared to previous semantics, the \semrule{Async-App} rule does a lot of heavy lifting. The newly created task frame is not immediately run, rather it is pushed onto the task queue $Q$ --- modeling semi-eager async function application.

There are two lines that compose Swifts enforcement of \dyn \extent. First, the \task allocated is a \emph{dependency} of the currently executing frame labeled $\ell$. Second, when the application is finished the new \task, $x_{\mathit{task}}$, awaits all of its dependencies before returning and rescheduling it's dependents. Furthermore, before awaiting its dependencies the \task will cancel it's dependencies as a way to model \cancelled \destruction. Although the dependencies are awaited, if an unawaited exception was raised, it will not propagate at the end of the scope. After awaiting dependencies the current \task immediately reschedules dependents to run.

We omit the \semrule{Await} rule as Swift uses \dyn \suspension and the rule is identical to the one presented for \CSharp.

The rules \semrule{Cancel} and \semrule{Cancelled?}, alongside the metafunction \code|task_cancelled| compose Swifts native cancellation mechanism. The key insight is that a cancellation check for a \task not only checks it's own cancellation state, but also that of it's dependents. When a \task is cancelled, all of its dependencies are simultaneously knowing of the cancellation request. Although Swift cancellation is cooperative, this built-in mechanism does not rely on \tasks to communicate cancellation to dependents/dependencies.

\subsection{Python}

\begin{figure}[H]
  \begin{subfigure}{\textwidth}
    \input{./f/py-gr.tex}
    \cprotect\caption{The py grammar extension on \exn.}
    \label{fig:py:grammar}
  \end{subfigure}\vspace{1em}
  \begin{subfigure}{\textwidth}
    \noindent
\begin{tabular*}{\linewidth}{@{} l @{\extracolsep{\fill}} r @{}}
$\langle \sigma_0, E[(\mathbf{async\ fn}\ \many{x} \to e_{\mathit{body}})\ \many{v}] \rangle \longrightarrow \langle \sigma_1, E[\ceag{\mathrm{reset}\ (\mathrm{shift}\ k \to k \mathbin{;} e_{\mathit{subst}})}] \rangle$ & [Async-App] \\
\multicolumn{2}{@{}l@{}}{
  $\begin{array}{l}
  \textbf{where}\\ 
  \qquad \many{x}_{\mathit{fresh}} = \many{\mathrm{gensym}()}\ \mathrm{and}\ \sigma_1 = \sigma_0[\many{x}_{\mathit{fresh}} \mapsto \many{v}] \\
  \qquad e_{\mathit{subst}} = e_{\mathit{body}}[\many{x}_{\mathit{fresh}}/\many{x}]
  \end{array}$
}
\end{tabular*}

\vspace{1.5em}

\noindent
\begin{tabular*}{\linewidth}{@{} l @{\extracolsep{\fill}} r @{}}
$\langle \sigma, E[\mathbf{await}\ (\mathrm{fun}\ x \to \mathrm{reset}\ (x \mathbin{;} e))] \rangle \longrightarrow \langle \sigma, E[e] \rangle$ & [Await-Coroutine]
\end{tabular*}

    \makelegend{eag}
    \cprotect\caption{}
    \label{fig:py:red}
  \end{subfigure}
  \cprotect\caption{The core operational semantics for py async/await.}
  \label{fig:lang:py}
\end{figure}

A core design choice of Python's async design is to separate the \coro mechanism --- generated by the compiler from an \psudo{async} function --- and the runtime. To that extent \Cref{fig:lang:py} provides the Python language extension, which provides the semantics for only an async application and awaiting a \coro.

The \semrule{Async-App} rule highlights the use of \code{shift} and \code{reset} to capture the function body as a continuation. In our model of the semantics we use the raw continuation as the underlying \coro value. Because the base language does not have a notion of \task, \coros can be only awaited to force a value. The \semrule{Await-Coroutine} rule removes the \code{reset} wrapper around the inner \coro to splice it's evaluation into the current context.

\paragraph{Asyncio}

\begin{figure}[H]
  \begin{subfigure}{\textwidth}
    $$
\begin{array}{rcl}
  e & ::= & \ldots \mid \mathbfsf{spawn}\ e \mid \mathbfsf{cancel}\ e
\end{array}
$$

    \cprotect\caption{The Asyncio grammar that extends the base Python language presented in \Cref{fig:py:grammar}. The \code{spawn} expression elevates a \coro into a \task and \code{cancel} marks a \task to be cancelled.}
    \label{fig:aio:grammar}
  \end{subfigure}\vspace{1em}
  \begin{subfigure}{\textwidth}
    \noindent
\begin{tabular*}{\linewidth}{@{} l @{\extracolsep{\fill}} r @{}}
$\langle t, \sigma_0, Q_0, T, P \uplus \{ \langle \ell, E[\mathrm{spawn}\ v_{\mathit{coro}}] \rangle \circ FS \} \rangle$ & [Spawn] \\
$\longrightarrow \langle \mathrm{step}(t), \sigma_1, Q_1, T, P \uplus \{ \langle \ell, E[x_{\mathit{task}}] \rangle \circ FS \} \rangle$ & \\
\multicolumn{2}{@{}l@{}}{
  $\begin{array}{l}
  \textbf{where}\ (\sigma_1, x_{\mathit{task}}) = \mathrm{task\text{-}alloc}(\sigma_0) \\
  \qquad \kappa = \mathbf{fn}\ x_{\mathit{unit}} \to x_{\mathit{unit}} \mathbin{;} \mathbfsf{reset}\ ( \\
  \qquad \qquad \mathbfsf{try}\ \mathrm{task\text{-}set\text{-}done}(x_{\mathit{task}}, \mathbf{await}\ v_{\mathit{coro}})\ \mathbfsf{catch}\ v_{\mathit{err}} \to \mathrm{task\text{-}set\text{-}failed}(x_{\mathit{task}}, v_{\mathit{err}}) \mathbin{;} \\
  \qquad \qquad \cext{\mathrm{os\text{-}start\text{-}soon}(\mathrm{task\text{-}get\text{-}dependents}(x_{\mathit{task}}))} \\
  \qquad ) \\
  \qquad Q_1 = Q_0 \qpush (x_{\mathit{task}}, \kappa)
  \end{array}$
}
\end{tabular*}

% \vspace{1.5em}
% \noindent
% \begin{tabular*}{\linewidth}{@{} l @{\extracolsep{\fill}} r @{}}
% $\langle t, \sigma, Q, T, P \uplus \{ \langle \ell, E[\mathbf{await}\ v_{\mathit{aw}}] \rangle \circ FS \} \rangle$ & [Await-Task] \\
% $\longrightarrow \langle \mathrm{step}(t), \sigma, Q, T, P \uplus \{ \langle \ell, E[e_{\mathit{check}}] \rangle \circ FS \} \rangle$ & \\
% \multicolumn{2}{@{}l@{}}{
%   $\begin{array}{l}
%   \textbf{where}\ \mathrm{task\text{-}is\text{-}task}(v_{\mathit{aw}}) \\
%   \qquad e_k = \mathbfsf{shift}\ k \to \mathrm{task\text{-}add\text{-}self\text{-}as\text{-}dependent}(v_{\mathit{aw}}, \ell, \mathrm{task\text{-}continue\text{-}with}(v_{\mathit{aw}}, k)) \\
%   \qquad e_{\mathit{check}} = \csus{\mathbfsf{if}\ \mathrm{task\text{-}is\text{-}completed}(v_{\mathit{aw}})\ \mathbfsf{then}\ \mathrm{task\text{-}get\text{-}result}(v_{\mathit{aw}})\ \mathbfsf{else}\ e_k}
%   % \mathrm{task\text{-}set\text{-}awaited}(v_{\mathit{aw}}) \mathbin{;} \\
%   \end{array}$
% }
% \end{tabular*}

\vspace{1.5em}
\noindent
\begin{tabular*}{\linewidth}{@{} l @{\extracolsep{\fill}} r @{}}
$\langle t, \sigma, (\ell_{\mathit{waiting}}, \kappa) \circ Q_1, T, P \uplus \{ \epsilon \} \rangle$ & [Schedule] \\
$\longrightarrow \langle \mathrm{step}(t), \sigma, Q_1, T, P \uplus \{ \langle \ell_{\mathit{waiting}}, \kappa\ () \rangle \circ \epsilon \} \rangle$ & \\
\multicolumn{2}{@{}l@{}}{
  $\begin{array}{l}
  \textbf{where}\ \neg \mathrm{task\text{-}cancelled}(\sigma, \ell_{\mathit{waiting}})
  \end{array}$
}
\end{tabular*}

\vspace{1.5em}
\noindent
\begin{tabular*}{\linewidth}{@{} l @{\extracolsep{\fill}} r @{}}
$\langle t, \sigma, (\ell_{\mathit{waiting}}, \kappa) \circ Q_1, T, P \uplus \{ \epsilon \} \rangle$ & [Schedule-Cancelled] \\
$\longrightarrow \langle \mathrm{step}(t), \sigma, Q_1, T, P \uplus \{ \langle \ell_{\mathit{waiting}}, \ccan{\mathbfsf{throw\text{-}in}(\kappa, \text{``cancelled''})} \rangle \circ \epsilon \} \rangle$ & \\
\multicolumn{2}{@{}l@{}}{
  $\begin{array}{l}
  \textbf{where}\ \mathrm{task\text{-}cancelled}(\sigma, \ell_{\mathit{waiting}})
  \end{array}$
}
\end{tabular*}

\vspace{1.5em}
\noindent
\begin{tabular*}{\linewidth}{@{} l @{\extracolsep{\fill}} r @{}}
$\langle t, \sigma \uplus \{ x \mapsto v \}, Q_{\mathit{before}} \circ (x, \kappa) \circ Q_{\mathit{after}}, T, P \rangle$ & [Sys-GC] \\
$\longrightarrow \langle \mathrm{step}(t), \sigma, Q_{\mathit{before}} \circ Q_{\mathit{after}}, T, P \rangle$ & \\
\multicolumn{2}{@{}l@{}}{
  $\begin{array}{l}
  \textbf{where}\ x \notin \crfe{\mathit{fv}(\sigma, T, P)}
  \end{array}$
}
\end{tabular*}

% \vspace{1.5em}
% \noindent
% \begin{tabular*}{\linewidth}{@{} l @{\extracolsep{\fill}} r @{}}
% $\langle t, \sigma, Q, T, P \uplus \{ \langle \ell, E[\mathrm{cancel}\ v_{\mathit{task}}] \rangle \circ FS \} \rangle$ & [Cancel] \\
% $\longrightarrow \langle t, \sigma, Q, T, P \uplus \{ \langle \ell, E[\mathrm{task\text{-}set\text{-}cancelled}(v_{\mathit{task}})] \rangle \circ FS \} \rangle$ & 
% \end{tabular*}

% \vspace{1.5em}
% \noindent
% \begin{tabular*}{\linewidth}{@{} l @{\extracolsep{\fill}} r @{}}
% $\langle t, \sigma, Q, T, P \uplus \{ \langle \mathrm{root}, E[\mathrm{os\text{-}block}\ v_{\mathit{coro}}] \rangle \circ FS \} \rangle$ & [OS-Block-Coro] \\
% $\longrightarrow \langle \mathrm{step}(t), \sigma, Q, T, P \uplus \{ \langle \mathrm{root}, E[\mathrm{os\text{-}block}\ (\mathrm{spawn}\ v_{\mathit{coro}})] \rangle \circ FS \} \rangle$ & \\
% \multicolumn{2}{@{}l@{}}{
%   $\begin{array}{l}
%   \textbf{where}\ v_{\mathit{coro}} = \mathrm{fun}\ x \to e
%   \end{array}$
% }
% \end{tabular*}

    \makelegend{eag,ext,ref,can}
    \cprotect\caption{}
    \label{fig:aio:red}
  \end{subfigure}
  \cprotect\caption{The core operational semantics for aio async/await.}
  \label{fig:lang:aio}
\end{figure}

\Cref{fig:lang:aio} provides the remaining grammar and operational semantics for Asyncio. The \semrule{Spawn} rule does not cancel/await dependencies when it finishes executing and instead immediately schedules dependents to run. We again have ommitted the \semrule{Await} rule as it is identical to that presented in \Cref{fig:cs:red}.

Rules \semrule{Schedule} and \semrule{Schedule-Cancelled} show how a \task is rescheduled onto an idle thread. The next available \task is taken from the task queue and scheduled on the next available thread. If the removed \task is cancelled, then it is not normally resumed, but rather it is resumed with an exception using the \code|throw_in| expression introduced in \exn.

\subsubsection{Trio}

\begin{figure}[H]
  \noindent
\begin{tabular*}{\linewidth}{@{} l @{\extracolsep{\fill}} r @{}}
$\langle t, \sigma_0, Q_0, T, P \uplus \{ \langle \ell, E[\mathrm{spawn}\ v_{\mathit{coro}}] \rangle \circ FS \} \rangle$ & [Spawn] \\
$\longrightarrow \langle \mathrm{step}(t), \sigma_1, Q_1, T, P \uplus \{ \langle \ell, E[x_{\mathit{task}}] \rangle \circ FS \} \rangle$ & \\
\multicolumn{2}{@{}l@{}}{
  $\begin{array}{l}
  \textbf{where}\ (\sigma_1, x_{\mathit{task}}) = \cext{\mathrm{task\text{-}alloc\text{-}dependency}(\sigma_0, \ell)} \\
  \qquad \kappa = \mathbf{fn}\ () \to \mathbfsf{reset}\ ( \\
  \qquad \qquad \mathbfsf{try}\ \mathrm{task\text{-}set\text{-}done}(x_{\mathit{task}}, \mathrm{await}\ v_{\mathit{coro}}) \mathbin{;} \cpro{\mathrm{task\text{-}await\text{-}dependencies}}\cext{(x_{\mathit{task}})} \\
  \qquad \qquad \mathbfsf{catch}\ v_{\mathit{err}} \to \mathrm{task\text{-}set\text{-}failed}(x_{\mathit{task}}, v_{\mathit{err}}) \mathbin{;} \\
  \qquad \qquad \mathrm{os\text{-}start\text{-}soon}(\mathrm{task\text{-}get\text{-}dependents}(x_{\mathit{task}})) \\
  \qquad ) \\
  \qquad Q_1 = Q_0 \qpush (x_{\mathit{task}}, \kappa)
  \end{array}$
}
\end{tabular*}

  \makelegend{ext,pro,can}
  \cprotect\caption{}
  \label{fig:trio:red}
\end{figure}
